\section*{Figure 1}
\textbf{Anatomical landmarks for all morphometric measures.} (A--G) Landmark pairs define the seven linear dimensions in Table~1. (H) The subpubic angle has vertex $p_2$ and arms toward $p_1$ and $p_3$; the view is normal to their plane. (I) Measurement definitions and landmark key. Markers and segments are overlaid on the reference pelvic geometry to keep internal landmarks visible. Values describe the template, not cohort averages. Lengths are measured in three dimensions; each camera direction is perpendicular to its measurement segment. Panels use different views and magnifications.

\section*{Figure 2}
\textbf{Anatomy, motion definitions and load application.} (A) Pelvic assembly: green denotes the superior S1 endplate constraint, orange SIJ cartilage and purple the symphysis. (B) Exploded articulation with illustrative axis directions; these arrows do not define the subject-specific PCA frame or anatomical signs. (C) Rotational and translational reporting conventions. (D,E) Bilateral and unilateral acetabular loads, each totaling 800~N. (F--H) Independent internal-ring, ischial and AP force pairs. Load arrows show applied global directions, not measured joint opening. The S1 penalty approximates a fixed boundary.

\section*{Figure 3}
\textbf{Predictive information in the anatomical triad.} (A,B) Held-out prediction $R^2$ for size and triad models, separately by load and endpoint. Each subject is held out from training in all loads. (C) Triad-model translations versus archived FE endpoints; each plotted prediction is averaged over validation repeats for display. The reported metrics in panels A,B and Table~3 pool squared errors over repeats rather than using these averaged predictions. The dashed line is identity. Predictor selection was retrospective; these are internal validation results.

\section*{Figure 4}
\textbf{Load-dependent motion and asymmetry.} Rotation-triplet norm, relative translation norm, absolute ML rotation and bilateral translation-magnitude asymmetry. Points are subjects; boxplots summarize each sex separately using medians and IQRs, with whiskers at 1.5 IQR. The numerical summaries in the text and Table~4 pool both sexes. Female and male groups use vermilion and blue. Jitter is deterministic. No signed nutation direction is inferred.

\section*{Figure 5}
\textbf{Paired absolute-component contrasts.} (A) Median within-subject differences from SP2leg, with bootstrap 95\% CIs. (B) Individual ML/CC contrast pairs and arrows to their componentwise medians. Positive values mean larger component magnitude, not outward joint motion. Corrected $q$ values are reported in Supplementary Table S1; intervals are pointwise, not simultaneous.

\section*{Figure 6}
\textbf{A directional limit to scalar response reduction.} Each 400~N transverse side force is progressively transferred from internal-ring to ischial patches, holding anatomy, tissues and total side force fixed. (A) Thin curves show subject translations re-extracted from exact mixtures of the FE displacement fields; thick curves show cohort medians for these translations and the two endpoint estimates. (B) Endpoint-vector alignment versus within-subject reduction at equal force redistribution. The cosine is averaged over the two sides; it is a directional diagnostic, not an independent predictor. (C) Scalar and signed-vector endpoint-combination errors. Field superposition is exact in the implemented linear system; combining extracted kinematic vectors is approximate, with its residual error measured.

\section*{Figure 7}
\textbf{Adjustment for volume and selected dimensions attenuates sex-associated translation differences.} (A,B) M2 female-minus-male coefficients with HC3 95\% CIs, shown on separate translation and rotation axes. Labels give BH-adjusted $q$ values across six tests. (C) Translation coefficients attenuate as pelvic volume and the selected triad enter the model; connecting lines compare specifications and are not causal pathways. (D) The positive pooled LAB1 subpubic-angle association is not reproduced within either sex. All M2 intervals span zero; this does not establish equivalence. M0 adjusts for age, M1 adds volume, and M2 adds the morphometric triad.

\section*{Figure 8}
\textbf{Retaining individual geometry preserves between-subject variation.} (A) Geometry-preserved / standardized-density and (B) density-preserved / template-geometry sample variances, each divided by the full-model variance, in percent. Five endpoints are shown. Color scales differ between panels and follow their actual data ranges. Values above 100\% are allowed. Subject-level agreement complements these variance summaries (translation RMSE $\leq0.0280$~mm; minimum identity-line $R^2=0.9457$). Bootstrap intervals and agreement diagnostics are provided in Supplementary Table S2; neither panel represents variance explained by an independent causal factor.
